% 1B mock exam (AP-style). Build: python3 units/ecology/study-guides/build.py 1B/exam
% Macros: latex/exam.sty. Figures: \drawfig (figures/ in this folder) and \bookcrop (cropped from a book by build.py).
\documentclass[11pt,letterpaper]{article}
\newcommand{\sgroot}{../}
% Shared preamble for the study guides and mock exams (built by ../build.py).
% \sgroot is the path from the document's folder to study-guides/ (empty for guides, ../ for exams).
\providecommand{\sgroot}{}
\usepackage[margin=0.8in,headheight=15pt]{geometry}
\usepackage{fontspec}
\setmainfont{texgyrepagella}[Extension=.otf,UprightFont=*-regular,BoldFont=*-bold,ItalicFont=*-italic,BoldItalicFont=*-bolditalic]
\setsansfont{texgyreheros}[Extension=.otf,UprightFont=*-regular,BoldFont=*-bold,ItalicFont=*-italic,BoldItalicFont=*-bolditalic]
\usepackage{amsmath,amssymb,booktabs,longtable,array,calc,enumitem,xcolor,graphicx}
\usepackage{tikz}
\usepackage{pgfplots}\pgfplotsset{compat=1.18}
\usetikzlibrary{calc,arrows.meta,positioning}
\usepackage[most]{tcolorbox}
\usepackage{fancyhdr,titlesec,needspace,etoolbox,environ}
\usepackage{hyperref}

\definecolor{forest}{HTML}{24594B}
\definecolor{pale}{HTML}{EFF5F1}
\definecolor{keyc}{HTML}{C98A0B}   % key idea: amber
\definecolor{vocabc}{HTML}{2F6FA8} % vocabulary: blue
\definecolor{tipc}{HTML}{3F8A3A}   % make it make sense: green
\definecolor{watchc}{HTML}{B8452F} % watch out: red
\definecolor{linkc}{HTML}{7B4C9A}  % connection: purple
\hypersetup{colorlinks=true,linkcolor=forest,urlcolor=vocabc}

\pagestyle{fancy}\fancyhf{}
\renewcommand{\headrulewidth}{0pt}
\fancyfoot[C]{\small\sffamily\thepage}
\setlength{\parindent}{0pt}\setlength{\parskip}{5pt}
\setlength{\emergencystretch}{2em}
\setlist[itemize]{leftmargin=1.3em,itemsep=2pt,topsep=2pt}
\setlist[enumerate]{leftmargin=1.6em,itemsep=2pt,topsep=2pt}
\renewcommand{\arraystretch}{1.2}
\providecommand{\tightlist}{\setlength{\itemsep}{0pt}\setlength{\parskip}{0pt}}
\providecommand{\pandocbounded}[1]{#1}
\def\LTcaptype{none}

\titleformat{\section}{\Large\sffamily\bfseries\color{forest}}{}{0pt}{}[{\color{forest}\titlerule[0.8pt]}]
\titleformat{\subsection}{\large\sffamily\bfseries\color{forest}}{}{0pt}{}
\titleformat{\subsubsection}{\normalsize\sffamily\bfseries}{}{0pt}{}
\titlespacing*{\section}{0pt}{14pt}{8pt}
\titlespacing*{\subsection}{0pt}{12pt}{4pt}
\setcounter{secnumdepth}{0}

% Block quotes (textbook excerpts) as shaded boxes.
\renewenvironment{quote}{\begin{tcolorbox}[breakable,enhanced,colback=pale,colframe=forest,boxrule=0pt,leftrule=2.5pt,arc=0pt,
  left=8pt,right=8pt,top=4pt,bottom=4pt,before skip=4pt,after skip=6pt]\small}{\end{tcolorbox}}

% ------------------------------------------------------------------ annotated textbook pages
% Landscape letter page: the scanned page on the left (full height), notes on the right.
\newlength{\portraitW}\setlength{\portraitW}{8.5in}
\newlength{\portraitH}\setlength{\portraitH}{11in}
\def\pgmargin{21.6}  % pt, top/left margin around the scan
\def\pgheight{568.8} % pt, printed height of the scan (8.5in - 2 margins)
\newcommand{\setpagesize}[2]{%
  \global\pdfpagewidth=#1\global\pdfpageheight=#2%
  \global\paperwidth=#1\global\paperheight=#2}

\newcommand{\circlednum}[2]{\tikz[baseline=(c.base)]\node[circle,fill=#1,text=white,font=\scriptsize\sffamily\bfseries,inner sep=1.2pt,minimum size=11pt](c){#2};}

% \textbookpage{image}{page width pt}{page height pt}{lesson}{book page}{intro}{highlights}{notes}
% Notes (\addnote) are typeset into boxes first, then laid out: each sits level with its highlight
% when there is room, is pushed down to avoid the note above, and is pulled back up from the page bottom.
\newcount\nnotes
\newbox\hdrbox
\newdimen\prevbottom \newdimen\limitdim \newdimen\tmpdim
\def\notegap{7pt}
\newcommand{\addnote}[5]{% color, number, label, anchor (scan pt, <0 = none), text
  \global\advance\nnotes by 1
  \expandafter\ifx\csname nb@\the\nnotes\endcsname\relax\expandafter\newbox\csname nb@\the\nnotes\endcsname\fi
  \global\expandafter\setbox\csname nb@\the\nnotes\endcsname=\hbox{\tikz\node[notebox=#1]{\notehead{#1}{#2}{#3}#5};}%
  \expandafter\xdef\csname na@\the\nnotes\endcsname{#4}%
  \expandafter\xdef\csname nc@\the\nnotes\endcsname{#1}}
\newcommand{\nb}[1]{\csname nb@#1\endcsname}
\newcommand{\layoutnotes}{%
  \global\prevbottom=\dimexpr\pgmargin pt+\ht\hdrbox+\dp\hdrbox\relax
  \ifnum\nnotes>0
    \foreach \k in {1,...,\the\nnotes}{%
      \edef\a{\csname na@\k\endcsname}%
      \tmpdim=\dimexpr\prevbottom+\notegap\relax
      \ifdim \a pt<0pt\else
        \pgfmathsetlength{\dimen0}{\pgmargin+\a*\s-7}%
        \ifdim\dimen0>\tmpdim \tmpdim=\dimen0 \fi
      \fi
      \expandafter\xdef\csname ny@\k\endcsname{\the\tmpdim}%
      \global\prevbottom=\dimexpr\tmpdim+\ht\nb{\k}+\dp\nb{\k}\relax}%
    \global\limitdim=\dimexpr612pt-32pt\relax
    \foreach \k in {\the\nnotes,...,1}{%
      \tmpdim=\csname ny@\k\endcsname
      \ifdim\dimexpr\tmpdim+\ht\nb{\k}+\dp\nb{\k}\relax>\limitdim
        \tmpdim=\dimexpr\limitdim-\ht\nb{\k}-\dp\nb{\k}\relax
        \expandafter\xdef\csname ny@\k\endcsname{\the\tmpdim}%
      \fi
      \global\limitdim=\dimexpr\tmpdim-\notegap\relax}%
    \ifdim\limitdim<\dimexpr\pgmargin pt+\ht\hdrbox+\dp\hdrbox-\notegap\relax
      \typeout{WARNING: notes overflow on textbook page \currentbookpage}\fi
  \fi}
\newcommand{\drawnotes}{%
  \ifnum\nnotes>0
    \foreach \k in {1,...,\the\nnotes}{%
      \edef\a{\csname na@\k\endcsname}\edef\c{\csname nc@\k\endcsname}\edef\y{\csname ny@\k\endcsname}%
      \node[anchor=north west,inner sep=0pt] at ($(O)+(\notex pt,-\y)$) {\usebox{\nb{\k}}};
      \ifdim \a pt<0pt\else
        \draw[\c,line width=0.6pt] ($(scan.north east)+(0,{-\a*\s pt})$) -- ++(6pt,0) -- ($(O)+(\notex pt-1pt,-\y-7pt)$);
      \fi}%
  \fi}
\newcommand{\textbookpage}[8]{%
  \clearpage\setpagesize{11in}{8.5in}\thispagestyle{empty}%
  \def\currentbookpage{#5}%
  \pgfmathsetmacro{\s}{\pgheight/#3}%
  \pgfmathsetmacro{\pgright}{\pgmargin+#2*\s}%
  \pgfmathsetmacro{\notex}{\pgright+20}%
  \pgfmathsetmacro{\notew}{792-\notex-\pgmargin}%
  \pgfmathsetmacro{\notetw}{\notew-10}%
  \setbox\hdrbox=\vbox{\hsize=\notew pt\parskip=0pt
    {\raggedright\sffamily\bfseries\color{forest}#4\par}\vspace{1pt}
    {\sffamily\footnotesize\color{black!55}Miller \& Levine Biology, textbook page #5\par}\vspace{-1pt}
    {\color{forest}\rule{\notew pt}{0.6pt}}\par\vspace{2pt}
    \ifx\relax#6\relax\else{\small\raggedright #6\par}\fi}%
  \global\nnotes=0 #8\layoutnotes
  \null
  \begin{tikzpicture}[remember picture,overlay]
    \coordinate (O) at (current page.north west);
    \node[anchor=north west,inner sep=0pt,draw=black!25] (scan) at ($(O)+(\pgmargin pt,-\pgmargin pt)$)
      {\includegraphics[height=\pgheight pt]{#1}};
    \node[anchor=north west,inner sep=0pt] at ($(O)+(\notex pt,-\pgmargin pt)$) {\usebox{\hdrbox}};
    \node[anchor=south east,font=\scriptsize\sffamily,text=black!50,inner sep=0pt] at ($(O)+(792pt-\pgmargin pt,-603pt)$)
      {Part 1 · the textbook, annotated \quad·\quad Miller \& Levine Biology (2019), p.\,#5 \quad·\quad page \thepage};
    #7
    \drawnotes
  \end{tikzpicture}%
  \clearpage\setpagesize{\portraitW}{\portraitH}}

% Highlight one text line on the scan: \hl{color}{x0}{y0}{x1}{y1} in scan points (y down).
\newcommand{\hl}[5]{%
  \fill[#1,opacity=0.22] ($(scan.north west)+({(#2-1.5)*\s pt},{-(#3-1)*\s pt})$) rectangle
                         ($(scan.north west)+({(#4+1.5)*\s pt},{-(#5+1)*\s pt})$);}
% Numbered badge just left of a highlight.
\newcommand{\badge}[4]{%
  \node[circle,fill=#1,text=white,font=\scriptsize\sffamily\bfseries,inner sep=1pt,minimum size=10pt]
    at ($(scan.north west)+({#3*\s pt-7pt},{-#4*\s pt})$) {#2};}
\newcommand{\notehead}[3]{{\sffamily\bfseries\footnotesize\circlednum{#1}{#2}\ \color{#1}\MakeUppercase{#3}}\\[1pt]}
\tikzset{notebox/.style={text width=\notetw pt,inner sep=4pt,fill=#1!7,draw=#1!55,line width=0.4pt,
  font=\small,align=left,rounded corners=1.5pt}}

% ------------------------------------------------------------------ figures in the guide text
% Caption used under both book figures and new diagrams.
\newcommand{\figcaption}[2]{\par\smallskip{\small\sffamily\raggedright #1\par}%
  \ifx\relax#2\relax\else{\scriptsize\sffamily\color{black!55}Source: #2\par}\fi}
% \bookfig{png}{width fraction}{caption}{source}: a figure cropped from one of the books.
\newcommand{\bookfig}[4]{\par\Needspace{8\baselineskip}\medskip
  \begin{minipage}{\linewidth}\centering
    \fcolorbox{black!15}{white}{\includegraphics[width=#2\linewidth]{#1}}
    \figcaption{#3}{#4}
  \end{minipage}\par\medskip}
% Several book figures side by side: \bookfigrow{\bookfigcol{w}{png}{height}{caption}{source}\hfill ...}
\newcommand{\bookfigcol}[5]{\begin{minipage}[t]{#1\linewidth}\centering
    \fcolorbox{black!15}{white}{\includegraphics[height=#3,width=0.95\linewidth,keepaspectratio]{#2}}
    \figcaption{#4}{#5}\end{minipage}}
\newcommand{\bookfigrow}[1]{\par\Needspace{10\baselineskip}\medskip\noindent#1\par\medskip}
% New diagrams drawn for the guide: \begin{guidefig}...\end{guidefig}{caption}
\newenvironment{guidefig}[1]{\par\Needspace{10\baselineskip}\medskip\begin{minipage}{\linewidth}\centering\def\guidefigcap{#1}}
  {\figcaption{\guidefigcap}{}\end{minipage}\par\medskip}
\tikzset{
  gbox/.style={draw=#1,fill=#1!8,rounded corners=3pt,align=center,font=\small\sffamily,inner sep=5pt},
  garr/.style={-{Stealth[length=2.4mm]},line width=1pt,draw=#1},
  glab/.style={font=\scriptsize\sffamily,fill=white,inner sep=1.5pt,align=center},
}

% ------------------------------------------------------------------ figures shared by guides and exams
% \drawfig{name}: a TikZ/pgfplots figure in figures/name.tex
\newcommand{\drawfig}[1]{\input{\sgroot figures/#1}}
% \bookcrop[width]{id}{book}{page}{x0 y0 x1 y1}{label}{caption}: a figure cropped from a book page.
% book: C = Campbell Biology, ML = Miller & Levine Biology; page = printed page number;
% crop box in points from the page's top-left corner. build.py makes build/figs/<id>.png from these arguments.
\newcommand{\bookname}[1]{\ifstrequal{#1}{C}{Campbell Biology}{Miller \& Levine Biology}}
\newcommand{\bookcrop}[7][0.7]{\bookfig{\sgroot build/figs/#2.png}{#1}{#7}{\bookname{#3}, #6, p.\,#4}}

\usepackage{../latex/exam}
\hypersetup{pdftitle={1B Mock Exam},pdfauthor={Prepared for Shawn}}

\begin{document}
\examtitle{1B}{1B Mock Exam: Characteristics of Life, Metabolism, Homeostasis, and Feedback}
\textbf{Format.} This practice exam follows the AP Biology exam format: data-based multiple choice, then free response. It covers learning target 1B: the characteristics of life, metabolism, homeostasis, negative and positive feedback, models, and controlled experiments.

\begin{center}
\begin{tabular}{@{}lll@{}}
\toprule
Section & Questions & Suggested time \\
\midrule
I. Multiple choice & 16 & 25 minutes \\
II. Free response & 4 (one long, three short) & 45 minutes \\
\bottomrule
\end{tabular}
\end{center}

Answer in full sentences in Section II. Graphs marked \textbf{hypothetical} were made up for practice. The answer key, with explanations and scoring guides, starts after Section II. Don't look until you've finished.

\clearpage

\examsection{Section I. Multiple choice}

Choose the \textbf{one} best answer for each question.

\questiongroup{Questions 1–3 refer to the characteristics of life.}

\begin{mcq}{1}
A student claims that fire is alive because it grows, uses oxygen, and releases energy. Which fact provides the \textbf{strongest} evidence against this claim?
\begin{choices}
  \item Fire gives off light.
  \item Fire is not made of cells and has no genetic code.
  \item Fire can move from place to place.
  \item Fire needs fuel to keep burning.
\end{choices}
\end{mcq}

\begin{mcq}{2}
A mule is the sterile offspring of a horse and a donkey. It cannot reproduce. Why do biologists still classify a mule as a living organism?
\begin{choices}
  \item Mules are large animals.
  \item The characteristics of life describe living things in general; a mule is made of cells, uses energy, maintains homeostasis, and has DNA.
  \item Any organism that can move is alive.
  \item Mules will evolve the ability to reproduce during their lifetimes.
\end{choices}
\end{mcq}

\begin{mcq}{3}
Which characteristic of life do viruses show \textbf{only} when they are inside a host cell?
\begin{choices}
  \item Being made of cells
  \item Maintaining homeostasis
  \item Reproducing
  \item Carrying out their own metabolism
\end{choices}
\end{mcq}

\questiongroup{Questions 4–5 refer to metabolism.}

\begin{mcq}{4}
Which process is an example of a metabolic reaction that \textbf{builds} large molecules and \textbf{uses} energy?
\begin{choices}
  \item Cellular respiration breaking glucose into carbon dioxide and water
  \item Digestion of starch into sugars
  \item Linking amino acids together to make a protein
  \item Breaking down glycogen to release glucose
\end{choices}
\end{mcq}

\begin{mcq}{5}
Yeast cells release carbon dioxide as they break down sugar during cellular respiration. Students measured the CO\textsubscript{2} released by yeast at different temperatures.
\drawfig{exam-yeast}

Which conclusion is best supported by the data?
\begin{choices}
  \item Yeast metabolism increases steadily with temperature up to 50\,°C.
  \item Metabolic activity was highest around 30–40\,°C and dropped sharply at 50\,°C.
  \item The difference between 30\,°C and 40\,°C shows that 40\,°C is clearly the best temperature.
  \item Temperature has no effect on yeast metabolism.
\end{choices}
\end{mcq}

\questiongroup{Questions 6–9 refer to the graph below.}

A runner's core body temperature and sweat rate were measured every 10 minutes. The runner exercised from minute 10 to minute 40.

\drawfig{exam-exercise}

\begin{mcq}{6}
Which statement best describes how the body responded to exercise?
\begin{choices}
  \item Body temperature rose without limit during exercise.
  \item Sweat rate increased after body temperature began to rise, and body temperature leveled off near 38.1\,°C.
  \item Sweat rate increased before exercise began.
  \item Body temperature fell during exercise because of sweating.
\end{choices}
\end{mcq}

\begin{mcq}{7}
The runner's temperature control is an example of which type of regulation?
\begin{choices}
  \item Positive feedback, because the sweat rate increased
  \item Negative feedback, because sweating opposes the rise in body temperature
  \item Positive feedback, because temperature stayed high
  \item No feedback, because temperature changed
\end{choices}
\end{mcq}

\begin{mcq}{8}
In this feedback loop, which part of the body acts as the \textbf{sensor and control center}?
\begin{choices}
  \item The sweat glands
  \item The skeletal muscles
  \item The hypothalamus in the brain
  \item The skin surface
\end{choices}
\end{mcq}

\begin{mcq}{9}
Why didn't the runner's temperature return to 37\,°C \textbf{immediately} when sweating increased?
\begin{choices}
  \item Physiological responses take time, so homeostasis reduces changes but doesn't eliminate them.
  \item Sweating raises body temperature.
  \item The set point changed to 38.1\,°C permanently.
  \item Feedback loops only work when the body is resting.
\end{choices}
\end{mcq}

\questiongroup{Questions 10–12 refer to the graph below.}

Two people each drank the same glucose solution after fasting overnight. Their blood glucose was measured every 30 minutes.

\drawfig{exam-glucose}

\begin{mcq}{10}
Which statement best describes Person X's results?
\begin{choices}
  \item Blood glucose rose after the drink, then returned to the normal range within about 90 minutes.
  \item Blood glucose stayed exactly the same during the whole test.
  \item Blood glucose kept rising for the whole 3 hours.
  \item Blood glucose fell below the normal range after 30 minutes.
\end{choices}
\end{mcq}

\begin{mcq}{11}
Which hypothesis best explains Person Y's results?
\begin{choices}
  \item Person Y's pancreas releases too much glucagon after meals.
  \item Person Y does not release enough insulin, or their cells do not respond to it, so glucose is not removed from the blood effectively.
  \item Person Y's negative feedback loop is working better than Person X's.
  \item Person Y drank more glucose than Person X.
\end{choices}
\end{mcq}

\begin{mcq}{12}
Between meals, blood glucose falls. Which response restores it?
\begin{choices}
  \item The pancreas releases insulin, and cells store more glucose.
  \item The pancreas releases glucagon, and the liver breaks down glycogen and releases glucose into the blood.
  \item The liver stores more glycogen.
  \item Muscles shiver to produce glucose.
\end{choices}
\end{mcq}

\questiongroup{Questions 13–16}

\begin{mcq}{13}
Which of the following is an example of \textbf{positive} feedback?
\begin{choices}
  \item Sweating when body temperature rises
  \item Releasing insulin when blood glucose rises
  \item During childbirth, contractions push the baby against the cervix, which triggers the release of more oxytocin and even stronger contractions
  \item Feeling thirsty when the body's water level drops
\end{choices}
\end{mcq}

\begin{mcq}{14}
Warming of the Arctic melts ice, exposing darker ocean water that absorbs more sunlight, which causes more warming. Which statement is correct?
\begin{choices}
  \item This is negative feedback, because it reduces the amount of ice.
  \item This is positive feedback, because the result reinforces the original change.
  \item This is homeostasis, because the ocean stays in balance.
  \item This is not a feedback loop, because it involves the environment instead of an organism.
\end{choices}
\end{mcq}

\begin{mcq}{15}
A student uses a home thermostat to explain homeostasis. Which statement is a \textbf{limitation} of this model?
\begin{choices}
  \item The thermostat has a sensor, a set point, and a response, like the body.
  \item The thermostat keeps temperature within a range.
  \item The body uses many sensors and several responses at once (sweating, shivering, changing blood flow), while a simple thermostat controls one heater.
  \item The thermostat uses negative feedback.
\end{choices}
\end{mcq}

\begin{mcq}{16}
A student tests whether a sports drink keeps body temperature lower during exercise than water does. Which design is the best controlled?
\begin{choices}
  \item Group 1 drinks the sports drink and runs outside on a hot day; Group 2 drinks water and walks indoors.
  \item Both groups do the same exercise in the same room for the same time; one group drinks the sports drink and the other drinks the same amount of water; temperature is measured the same way in both.
  \item One student drinks the sports drink and another drinks water; the student who feels better wins.
  \item Both groups drink the sports drink, and temperature is measured before and after.
\end{choices}
\end{mcq}

\clearpage

\examsection{Section II. Free response}

\frq{Question 1 (long response, 10 points)}

Students test how temperature affects the metabolism of yeast. They add the same mass of yeast and the same volume of sugar water to 15 flasks, place 5 flasks in each of three water baths (20\,°C, 30\,°C, and 40\,°C), and measure the CO\textsubscript{2} released every 5 minutes.

\begin{center}
\begin{tabular}{@{}lccccc@{}}
\toprule
Time (min) & 0 & 5 & 10 & 15 & 20 \\
\midrule
Mean CO\textsubscript{2} at 20\,°C (mL) & 0 & 1.2 & 2.6 & 4.1 & 5.8 \\
Mean CO\textsubscript{2} at 30\,°C (mL) & 0 & 2.9 & 6.1 & 9.0 & 12.2 \\
Mean CO\textsubscript{2} at 40\,°C (mL) & 0 & 3.1 & 6.4 & 9.6 & 12.9 \\
\bottomrule
\end{tabular}
\end{center}

\datanote{Hypothetical data. n = 5 flasks per temperature.}

\begin{parts}
  \item \textbf{Identify} the independent variable, the dependent variable, and one controlled variable.
  \item \textbf{Explain} why CO\textsubscript{2} production can be used as a measure of the yeast's metabolism.
  \item \textbf{Construct} a line graph of the data on the grid below, with one line for each temperature. Label the axes and include a key.

  \graphgrid
  \item \textbf{Calculate} the mean rate of CO\textsubscript{2} production (mL/min) at 20\,°C and at 30\,°C over the 20 minutes.
  \item \textbf{Describe} the effect of temperature on the rate of metabolism, using your answer to (d).
  \item A student predicts that at 60\,°C the yeast will release even more CO\textsubscript{2}. \textbf{Evaluate} this prediction, using what you know about temperature and living cells.
  \item \textbf{Describe} a control that would show that the CO\textsubscript{2} comes from the yeast and not from the sugar water itself.
\end{parts}

\frq{Question 2 (short response, 4 points)}

A student finds an unknown object in a pond sample. Under the microscope it has a cell membrane and moves. Over several days, the number of these objects in the sample increases.

\begin{parts}
  \item \textbf{Identify} two characteristics of life the object shows.
  \item \textbf{Describe} one additional observation or test that would give stronger evidence that the object is alive.
  \item \textbf{Explain} why one characteristic alone is not enough to decide whether something is alive.
  \item \textbf{Explain} why viruses are often described as being ``at the border'' of life.
\end{parts}

\frq{Question 3 (short response, 4 points)}

Refer to the diagram of blood glucose regulation.

\bookcrop[0.62]{c-glucose-exam}{C}{928}{190 400 614 754}{Figure 41.23}
  {Regulation of blood glucose by the pancreas, liver, and body cells.}

\begin{parts}
  \item \textbf{Identify} the stimulus that causes the pancreas to release insulin.
  \item \textbf{Describe} how insulin returns blood glucose to the normal range.
  \item \textbf{Explain} why this system is an example of negative feedback.
  \item \textbf{Predict} what would happen to blood glucose after a meal in a person whose body cells stopped responding to insulin, and \textbf{justify} your prediction.
\end{parts}

\frq{Question 4 (short response, 4 points)}

\begin{parts}
  \item \textbf{Define} homeostasis.
  \item Homeostasis is described as a \textbf{dynamic equilibrium}. \textbf{Explain} what this means, using body temperature as an example.
  \item \textbf{Contrast} the effect of negative feedback and positive feedback on a system.
  \item \textbf{Explain} why positive feedback is useful during childbirth even though it does not maintain homeostasis.
\end{parts}

\answerkey

\keyheading{Section I at a glance}

\begin{center}
\begin{tabular}{@{}lccccccc@{}}
\toprule
Q & Answer & Q & Answer & Q & Answer & Q & Answer \\
\midrule
1 & B & 5 & B & 9 & A & 13 & C \\
2 & B & 6 & B & 10 & A & 14 & B \\
3 & C & 7 & B & 11 & B & 15 & C \\
4 & C & 8 & C & 12 & B & 16 & B \\
\bottomrule
\end{tabular}
\end{center}

\textbf{Scoring guide.} 14–16 correct: ready. 11–13: review the explanations for the ones you missed. 10 or fewer: reread the study guide's Parts 2–3, then retake.

\keyheading{Section I explanations}

\answer{1}{B}{}
\skill{apply the criteria for life (M\&L p. 22).}
Living things are made of \textbf{cells} and have a \textbf{universal genetic code}. Fire has neither. A single missing criterion rules something out, so (B) is the strongest evidence.
\begin{whynot}
  \wrong{(A), (C), (D)}{these describe things fire \emph{does}, so they don't argue against the claim. Some living things also give off light, move, or need fuel.}
\end{whynot}

\answer{2}{B}{}
\skill{the criteria describe life in general.}
A mule meets every other criterion, and it came from parents that reproduced. The criteria describe living things as a group, not every individual.
\begin{whynot}
  \wrong{(A)}{size isn't a criterion.}
  \wrong{(C)}{movement alone isn't enough (cars and clouds move).}
  \wrong{(D)}{individuals don't evolve; groups do, over generations.}
\end{whynot}

\answer{3}{C}{Reproducing}
\skill{why viruses are borderline.}
Viruses ``can reproduce'' only ``after infecting living things'' (M\&L p. 24): they use the host cell's machinery.
\begin{whynot}
  \wrong{(A), (B), (D)}{viruses never have cells, homeostasis, or their own metabolism, even inside a host.}
\end{whynot}

\answer{4}{C}{}
\skill{two directions of metabolism (Campbell p. 146).}
Joining amino acids into a protein \textbf{builds} a large molecule and \textbf{uses} energy: an anabolic (building-up) reaction.
\begin{whynot}
  \wrong{(A), (B), (D)}{all break large molecules into smaller ones and release energy (catabolic).}
\end{whynot}

\answer{5}{B}{}
\skill{interpret a graph with error bars.}
CO\textsubscript{2} rose from 10 to 30\,°C, was similar at 30 and 40\,°C (the error bars overlap), and \textbf{crashed at 50\,°C}. High heat damages the enzymes that carry out metabolism.
\begin{whynot}
  \wrong{(A)}{the 50\,°C bar is lowest, not highest.}
  \wrong{(C)}{the 30 and 40\,°C error bars overlap, so the data don't show a clear difference.}
  \wrong{(D)}{there is a large effect.}
\end{whynot}

\answer{6}{B}{}
\skill{describe a trend in two variables.}
Temperature rose from 37.1\,°C at 10 minutes, the sweat rate climbed at 20–30 minutes, and temperature \textbf{leveled off} at 38.1\,°C instead of continuing up. After exercise both returned toward normal.
\begin{whynot}
  \wrong{(A)}{temperature leveled off.}
  \wrong{(C)}{sweat rate was low (0.1–0.2 L/h) before exercise.}
  \wrong{(D)}{temperature rose during exercise, although sweating limited the rise.}
\end{whynot}

\answer{7}{B}{}
\skill{identify negative feedback.}
The response (sweating) \textbf{opposes} the change (rising temperature). That is negative feedback, ``a stimulus produces a response that opposes the original stimulus'' (M\&L p. 907).
\begin{whynot}
  \wrong{(A), (C)}{positive feedback would make temperature rise faster and faster.}
  \wrong{(D)}{a changing value doesn't mean there's no feedback. Feedback responds \emph{to} change.}
\end{whynot}

\answer{8}{C}{The hypothalamus}
\skill{parts of a feedback loop.}
``A part of the brain called the hypothalamus contains nerve cells that monitor body temperature'' (M\&L p. 908), and it sends the signals that start sweating.
\begin{whynot}
  \wrong{(A)}{sweat glands are \textbf{effectors}: they carry out the response.}
  \wrong{(B)}{skeletal muscles are effectors for shivering.}
  \wrong{(D)}{the skin is where cooling happens, not where the decision is made.}
\end{whynot}

\answer{9}{A}{}
\skill{homeostasis as a dynamic equilibrium.}
``Physiological responses to stimuli are not instantaneous \dots{} homeostasis moderates but doesn't eliminate changes'' (Campbell p. 894).
\begin{whynot}
  \wrong{(B)}{sweating cools the body.}
  \wrong{(C)}{temperature returned to 37.1\,°C after exercise, so the set point didn't change permanently.}
  \wrong{(D)}{feedback is most important during challenges like exercise.}
\end{whynot}

\answer{10}{A}{}
\skill{read a line graph.}
Person X peaked at 150 mg/100 mL at 30 minutes and was back to 100 (inside the normal band) by 90 minutes.
\begin{whynot}
  \wrong{(B)}{glucose clearly changed.}
  \wrong{(C)}{it fell after 30 minutes.}
  \wrong{(D)}{it stayed at or above 84, inside the normal range.}
\end{whynot}

\answer{11}{B}{}
\skill{propose an explanation from data.}
Person Y started high (120), rose higher (235), and stayed far above the normal range for 3 hours. Glucose isn't being removed, which is what insulin does. This pattern is typical of diabetes.
\begin{whynot}
  \wrong{(A)}{extra glucagon would raise glucose, but the main problem shown is failure to \emph{lower} it after the drink; (B) explains that more directly.}
  \wrong{(C)}{Person Y's regulation is clearly weaker.}
  \wrong{(D)}{the question says both drank the same solution.}
\end{whynot}

\answer{12}{B}{}
\skill{the opposing hormone.}
When glucose drops, ``the pancreas releases glucagon. Glucagon stimulates liver and skeletal muscle cells to break down glycogen and release glucose into the blood'' (M\&L p. 932).
\begin{whynot}
  \wrong{(A), (C)}{these lower blood glucose.}
  \wrong{(D)}{shivering produces heat, not glucose.}
\end{whynot}

\answer{13}{C}{}
\skill{identify positive feedback.}
In childbirth, each contraction causes more oxytocin, which causes stronger contractions. The response \textbf{amplifies} the stimulus (Campbell p. 894).
\begin{whynot}
  \wrong{(A), (B), (D)}{each response opposes a change, which is negative feedback.}
\end{whynot}

\answer{14}{B}{}
\skill{feedback in a non-living system (teacher's slide 7).}
Warming $\rightarrow$ less ice $\rightarrow$ darker surface $\rightarrow$ more heat absorbed $\rightarrow$ more warming. The result \textbf{reinforces} the original change.
\begin{whynot}
  \wrong{(A)}{the loop amplifies warming; it doesn't oppose it.}
  \wrong{(C)}{nothing is being kept stable.}
  \wrong{(D)}{feedback loops occur in any system, including Earth's climate.}
\end{whynot}

\answer{15}{C}{}
\skill{evaluate a model.}
Every model leaves something out. The body uses many sensors and several responses (sweating, shivering, changing blood flow to the skin), while a simple thermostat switches one heater.
\begin{whynot}
  \wrong{(A), (B), (D)}{these are \emph{strengths} of the model: ways it matches the body.}
\end{whynot}

\answer{16}{B}{}
\skill{design a controlled experiment.}
In (B) the only difference between groups is the drink (the independent variable); exercise, room, time, amount drunk, and measuring method are all controlled.
\begin{whynot}
  \wrong{(A)}{three things differ (drink, temperature, and type of exercise).}
  \wrong{(C)}{one person per group is too small a sample, and ``feels better'' isn't measurable.}
  \wrong{(D)}{no comparison group.}
\end{whynot}

\keyheading{Section II scoring guides and model answers}

\scoreheading{Question 1 (10 points)}

\begin{center}
\begin{tabularx}{\linewidth}{@{}llY@{}}
\toprule
Part & Points & Earn the point(s) for \dots{} \\
\midrule
(a) & 1 & Independent: temperature. Dependent: volume of CO\textsubscript{2} released. Controlled: any one (mass of yeast, volume or concentration of sugar water, flask size, time). All three needed for the point. \\
(b) & 1 & Yeast release CO\textsubscript{2} when they break down sugar in cellular respiration, a metabolic reaction, so more CO\textsubscript{2} means faster metabolism. \\
(c) & 2 & Correct labeled axes with units: x = time (min), y = CO\textsubscript{2} released (mL) (1). Three correctly plotted lines with a key (1). \\
(d) & 2 & 20\,°C: 5.8 ÷ 20 = \textbf{0.29 mL/min} (1). 30\,°C: 12.2 ÷ 20 = \textbf{0.61 mL/min} (1). \\
(e) & 1 & Raising the temperature from 20\,°C to 30\,°C about doubled the rate; 40\,°C (0.65 mL/min) was only slightly faster than 30\,°C. \\
(f) & 2 & The prediction is not supported (1). Reasoning: very high temperatures damage the enzymes and other proteins that carry out metabolism, so the rate would likely drop sharply, as it did at 50\,°C in the multiple-choice data (1). \\
(g) & 1 & Flasks with sugar water but no yeast, kept at each temperature. They should release no CO\textsubscript{2}. \\
\bottomrule
\end{tabularx}
\end{center}

\modelanswer
\begin{parts}
  \item The independent variable is the temperature of the water bath. The dependent variable is the volume of CO\textsubscript{2} released. The mass of yeast should be kept the same in every flask.
  \item Yeast break down sugar in cellular respiration and release CO\textsubscript{2} as a product. Cellular respiration is part of metabolism, so the faster CO\textsubscript{2} is released, the faster the yeast's metabolism.
  \item \emph{A line graph: x-axis ``Time (min)'' from 0 to 20; y-axis ``CO\textsubscript{2} released (mL)'' from 0 to at least 13; three lines (20, 30, 40\,°C) with a key. The 30 and 40\,°C lines are close together and much steeper than the 20\,°C line.}
  \item 20\,°C: 5.8 mL ÷ 20 min = 0.29 mL/min. 30\,°C: 12.2 mL ÷ 20 min = 0.61 mL/min.
  \item Warmer temperatures sped up metabolism: going from 20\,°C to 30\,°C roughly doubled the rate (0.29 $\rightarrow$ 0.61 mL/min). Going from 30\,°C to 40\,°C changed it only a little (0.61 $\rightarrow$ 0.65 mL/min), so the rate levels off.
  \item I disagree. Metabolism depends on enzymes, which are proteins that work best within a certain temperature range. At 60\,°C the heat would likely damage the enzymes (and kill the yeast), so CO\textsubscript{2} release would drop sharply instead of increasing.
  \item Set up flasks with sugar water but no yeast at each temperature. If they release no CO\textsubscript{2}, the CO\textsubscript{2} in the other flasks must come from the yeast.
\end{parts}

\scoreheading{Question 2 (4 points)}

\begin{center}
\begin{tabularx}{\linewidth}{@{}lY@{}}
\toprule
Part & Earn 1 point for \dots{} \\
\midrule
(a) & Two of: made of a cell (has a cell membrane); responds to the environment or moves; reproduces (numbers increase). \\
(b) & A valid test for another criterion, e.g., check for DNA, measure use of food or oxygen (metabolism), show a response to a stimulus such as light, or watch it grow. \\
(c) & Nonliving things can share single traits with living things (fire grows and uses energy; cars move), so all the criteria are needed together. \\
(d) & Viruses reproduce and evolve (and have genetic material) but only inside a host cell; they have no cells and no metabolism of their own. \\
\bottomrule
\end{tabularx}
\end{center}

\modelanswer
\begin{parts}
  \item It seems to be made of a cell (it has a cell membrane), and it reproduces (the number increases over several days). Its movement may also be a response to its environment.
  \item I could test whether it contains DNA, or put a light at one side of the sample to see if the objects move toward or away from it (a response to a stimulus). Seeing each object take in food and grow would show metabolism and growth.
  \item Many nonliving things have one or two traits of life: fire grows and uses energy, and a car moves and responds to its driver. Only living things have all of the characteristics together, so one trait alone can't prove something is alive.
  \item Viruses have genetic material and can reproduce and evolve, but only after infecting a living cell and using its machinery. They are not made of cells and have no metabolism of their own, so they meet some criteria but not others.
\end{parts}

\scoreheading{Question 3 (4 points)}

\begin{center}
\begin{tabularx}{\linewidth}{@{}lY@{}}
\toprule
Part & Earn 1 point for \dots{} \\
\midrule
(a) & A rise in blood glucose (for example, after eating). \\
(b) & Insulin causes body cells to take up glucose and the liver to store it as glycogen, lowering blood glucose. \\
(c) & The response (lowering glucose) opposes the stimulus (high glucose), bringing the level back toward the set point. \\
(d) & Blood glucose would rise after a meal and stay high, because cells wouldn't take up glucose even with insulin present. \\
\bottomrule
\end{tabularx}
\end{center}

\modelanswer
\begin{parts}
  \item An increase in blood glucose, such as after eating a meal.
  \item Insulin travels in the blood and signals body cells to take up glucose and the liver to store it as glycogen. Glucose leaves the blood, so the level falls back to the normal range (70–110 mg/100 mL).
  \item The stimulus is high blood glucose, and the response (insulin causing glucose uptake) lowers it. The response opposes the change, so it is negative feedback. Once glucose is normal, insulin release decreases.
  \item Blood glucose would rise after the meal and stay high for a long time (like Person Y in Questions 10–11). Even if the pancreas released insulin, the cells wouldn't respond by taking up glucose, so glucose would stay in the blood.
\end{parts}

\scoreheading{Question 4 (4 points)}

\begin{center}
\begin{tabularx}{\linewidth}{@{}lY@{}}
\toprule
Part & Earn 1 point for \dots{} \\
\midrule
(a) & Keeping internal conditions relatively constant (within a range) despite changes in the internal or external environment. \\
(b) & Conditions constantly change a little and are constantly corrected, so they stay within a range; a correct body-temperature example. \\
(c) & Negative feedback opposes a change and stabilizes the system; positive feedback amplifies a change and pushes the system further from its starting point. \\
(d) & Childbirth is a process that must be completed; positive feedback builds contractions until the baby is born, then the loop stops. \\
\bottomrule
\end{tabularx}
\end{center}

\modelanswer
\begin{parts}
  \item Homeostasis is the maintenance of relatively constant internal conditions, such as temperature, water, and blood glucose, within a normal range, even when the internal or external environment changes.
  \item Dynamic equilibrium means the system is balanced overall but always moving and adjusting. Body temperature rises a little during exercise and falls a little on a cold day, but sweating and shivering keep correcting it, so it stays around 36–38\,°C.
  \item Negative feedback responds to a change by pushing the condition back the opposite way, which keeps the system stable. Positive feedback responds to a change by pushing the condition further in the same direction, so the change gets bigger and bigger.
  \item Childbirth needs to \emph{finish} rather than stay stable. Each contraction triggers the release of more oxytocin, which causes stronger contractions, until the baby is born. Then the stimulus is gone and the loop ends. Positive feedback drives the process to completion quickly.
\end{parts}

\end{document}
